\ifdefined\gkver\else\def\gkver{student}\fi
\documentclass[\gkver]{gaokaozhenti}
\begin{document}

\chapter{2006年上海理综}
%% paperType: 理综物理部分

\begin{enumerate}
\item
%% number: 7
%% paperName: 2006年上海理综第 7 题 3 分
%% typeId: 1
%% type: 单选题
%% chapter: 直线运动
%% point: 运动图象
%% method: 无
%% score: 3
%% degree: 850
%% duplicateId: 0
%% body:
右图是一娱乐场的喷水滑梯。若忽略摩擦力，人从滑梯顶端滑下直到入水前，速度大小随时间变化的关系最接近图\xzanswer{B}

\onepicture[w=3.2cm]{figs/07.png}

\fourchoices[answer=B, ispicture=true, h=1.5cm]
{figs/07a.png}
{figs/07b.png}
{figs/07c.png}
{figs/07d.png}

%% answer: B
%% memo:
\memoanswer{
本题原卷及材料中未提供详解，待补充。
}

\item
%% number: 8
%% paperName: 2006年上海理综第 8 题 3 分
%% typeId: 1
%% type: 单选题
%% chapter: 曲线运动
%% point: 圆周运动的应用
%% method: 无
%% score: 3
%% degree: 650
%% duplicateId: 0
%% body:
一艘宇宙飞船在预定轨道上做匀速圆周运动，在该飞船的密封舱内，下列实验能够进行的是\xzanswer{C}

\fourpicture[showlabel=false, h=2.2cm]
{figs/08a.png}
{figs/08b.png}
{figs/08c.png}
{figs/08d.png}

\fourchoices[answer=C]
{宇航员生活废水过滤处理实验}
{研究动能与重力势能转化规律实验}
{探究感应电流的产生条件实验}
{血浆与血细胞自然分层实验}

%% answer: C
%% memo:
\memoanswer{
本题原卷及材料中未提供详解，待补充。
}

\item
%% number: 9
%% paperName: 2006年上海理综第 9 题 3 分
%% typeId: 1
%% type: 单选题
%% chapter: 电路
%% point: 电功电功率
%% method: 无
%% score: 3
%% degree: 650
%% duplicateId: 0
%% body:
夏天空调器正常工作时，制冷状态与送风状态交替运行。一空调器在不同工作状态下电功率随时间变化的关系见右图，此空调器运转 1 小时用电\xzanswer{B}

\onepicture[w=6.5cm]{figs/09.png}

\fourchoices[answer=B]
{$1.0\UkWh$}
{$1.5\UkWh$}
{$2.0\UkWh$}
{$2.5\UkWh$}

%% answer: B
%% memo:
\memoanswer{
本题原卷及材料中未提供详解，待补充。
}

\item
%% number: 10
%% paperName: 2006年上海理综第 10 题 3 分
%% typeId: 1
%% type: 单选题
%% chapter: 直线运动
%% point: 平均速度和瞬时速度
%% method: 无
%% score: 3
%% degree: 650
%% duplicateId: 0
%% body:
客车运能是指一辆客车单位时间最多能够运送的人数。某景区客运索道的客车容量为 50 人/车，它从起始站运行至终点站（右图）单程用时 10 分钟。该客车运行的平均速度和每小时的运能约为\xzanswer{A}

\onepicture[w=7.2cm]{figs/10.png}

\fourchoices[answer=A]
{$5\Ums$，300 人}
{$5\Ums$，600 人}
{$3\Ums$，300 人}
{$3\Ums$，600 人}

%% answer: A
%% memo:
\memoanswer{
本题原卷及材料中未提供详解，待补充。
}

\item
%% number: 11
%% paperName: 2006年上海理综第 11 题 3 分
%% typeId: 1
%% type: 单选题
%% chapter: 曲线运动
%% point: 线速度角速度
%% method: 无
%% score: 3
%% degree: 550
%% duplicateId: 0
%% body:
若以某固定点为起点画出若干矢量，分别代表运动质点在不同时刻的速度，则这些矢量的末端所形成的轨迹被定义为“速矢端迹”。由此可知\xzanswer{B}

①匀速直线运动的速矢端迹是线段

②匀加速直线运动的速矢端迹是射线

③匀速圆周运动的速矢端迹是圆

④简谐运动的速矢端迹是点

\fourchoices[answer=B]
{①②}
{②③}
{③④}
{①④}

%% answer: B
%% memo:
\memoanswer{
本题原卷及材料中未提供详解，待补充。
}

\item
%% number: 12
%% paperName: 2006年上海理综第 12 题 3 分
%% typeId: 1
%% type: 单选题
%% chapter: 曲线运动
%% point: 竖直平面圆周运动
%% method: 无
%% score: 3
%% degree: 550
%% duplicateId: 0
%% body:
“时空之旅”飞车表演时，演员驾着摩托车，在球形金属网内壁上下盘旋，令人惊叹不已。摩托车沿图示竖直轨道做圆周运动过程中\xzanswer{D}

\onepicture[w=4.0cm]{figs/12.png}

\fourchoices[answer=D]
{机械能一定守恒}
{其输出功率始终保持恒定}
{经过最低点时的向心力仅由支持力提供}
{通过最高点时的最小速度与球形金属网直径有关}

%% answer: D
%% memo:
\memoanswer{
本题原卷及材料中未提供详解，待补充。
}

\item
%% number: 19
%% paperName: 2006年上海理综第 19 题 3 分
%% typeId: 1
%% type: 单选题
%% chapter: 电路
%% point: 电路中的图象
%% method: 无
%% score: 3
%% degree: 850
%% duplicateId: 0
%% body:
右图中①～④是四个物体在通电后电流随电压变化的关系图，其中电流与电压之间关系符合欧姆定律的是\xzanswer{C}

\onepicture[w=7.0cm]{figs/19.png}

\fourchoices[answer=C]
{①}
{②}
{③}
{④}

%% answer: C
%% memo:
\memoanswer{
本题原卷及材料中未提供详解，待补充。
}

\item
%% number: 31
%% paperName: 2006年上海理综第 31 题 6 分
%% typeId: 3
%% type: 填空题
%% chapter: 电路
%% point: 串并联电路
%% method: 无
%% score: 6
%% degree: 650
%% duplicateId: 0
%% body:
研究性学习小组经讨论得出家庭用电不安全事例主要有：

电器漏电、线路老化、\tkanswer[1.8cm]{无接地线}、\tkanswer[3.6cm]{使用劣质开关或插座等}等。并提出猜想：[1] 家庭电路都连接有熔丝（保险丝），应该是安全的；[2] 熔丝只能保护电路，而不能保护人。某同学说：我家还安装了漏电保护器，很安全的。于是，大家决定探究熔丝和漏电保护器在电路中的作用，小组同学汇总资料并绘制电器漏电示意图如下：

\textbf{资料}

①熔丝是一种合金丝，当电路中的电流超过熔丝的额定电流时，它就会熔断。

②当通过人体的电流大于 $30\UmA$ 时，将有生命危险。

③漏电保护器是防止电流泄漏的装置，如果火线与零线中的电流不同，它将切断电源。

结合上述图、表，能够支持猜想 [2] 的资料有 \tkanswer[1.2cm]{①、②}（从①、②、③中选择），理由是 \tkanswer[6cm]{流过熔丝的电流未超过其额定电流，但流过人体的电流已超过 $30\UmA$}。

电路中，若用铜丝代替熔丝，可能出现的危害是 \tkanswer[6cm]{电流过大而不熔断，容易发生火灾事故}。

\onepicture[w=9.5cm, align=right]{figs/31.png}

%% answer: 
\jdanswer{
无接地线，使用劣质开关或插座等；①、②，流过熔丝的电流未超过其额定电流，但流过人体的电流已超过 $30\UmA$；电流过大而不熔断，容易发生火灾事故。
}

%% memo:
\memoanswer{
本题原卷及材料中未提供详解，待补充。
}

\item
%% number: 32
%% paperName: 2006年上海理综第 32 题 4 分
%% typeId: 3
%% type: 填空题
%% chapter: 电路
%% point: 串并联电路
%% method: 无
%% score: 4
%% degree: 650
%% duplicateId: 0
%% body:
结合 31 题的资料③，你认为下图中四种（A、B、C、D）漏电保护器的安装方法正确的是（　　）

细心的同学发现上图中有一处用电不安全隐患。请将图中的不安全因素用笔圈出。

\onepicture[w=8.5cm, align=right]{figs/32.png}

%% answer: 
\jdanswer{
B，地线接在自来水管上
}

%% memo:
\memoanswer{
本题原卷及材料中未提供详解，待补充。
}

\item
%% number: 34
%% paperName: 2006年上海理综第 34 题 5 分
%% typeId: 3
%% type: 填空题
%% chapter: 电路
%% point: 串并联电路
%% method: 无
%% score: 5
%% degree: 650
%% duplicateId: 0
%% body:
观察下图 A，若淋浴器发生漏电可导致其外壳带电，请分析其危险性：\tkanswer[5cm]{容易发生人体触电事故}。为了更加合理和安全地使用电热淋浴器，该研究小组在理论上设计了 B、C 两方案（见下图），图中的 $R$ 相当于保护电阻。根据已有知识判断，可行方案是 \tkanswer[1.2cm]{C}（填“B”，或“C”），理由是 \tkanswer[6cm]{发生漏电时人体分得的电压只有 $12\UV$，不会造成危险}。

\onepicture[w=10cm, align=right]{figs/34.png}

%% answer: 
\jdanswer{
容易发生人体触电事故；C，发生漏电时人体分得的电压只有 $12\UV$，不会造成危险。
}

%% memo:
\memoanswer{
本题原卷及材料中未提供详解，待补充。
}

\item
%% number: 51
%% paperName: 2006年上海理综第 51 题 3 分
%% typeId: 3
%% type: 填空题
%% chapter: 功和能
%% point: 动能定理
%% method: 建立模型
%% score: 3
%% degree: 450
%% duplicateId: 0
%% body:
关注身体健康，提高生活质量。体检对预防疾病、保护健康有重要作用。当某些体检指标偏离常态时，表明受检个体的身体状况可能出现异常。如：血红蛋白含量下降则会导致贫血；嗜中性粒细胞增多，说明身体有炎症；嗜酸性粒细胞增多常见于一些过敏性疾病；血糖浓度偏高很可能患有糖尿病。

下面是某病人体检结果部分记录表，请分析回答问题：

\begin{center}
\begin{tblr}{|l|l|l|}
\hline
项目 & 测定值 & 正常范围 \\
\hline
红细胞个数 & 413 万个/立方毫米 & 400～500 个/立方毫米 \\
血红蛋白 & 0.1 毫克/立方毫米 & 0.12～0.16 毫克/立方毫米 \\
白细胞计数 & 11300 个/立方毫米 & 4000～10000 个/立方毫米 \\
嗜中性粒细胞绝对数 & 8300 个/立方毫米 & 2000～7500 个/立方毫米 \\
嗜酸性粒细胞绝对数 & 700 个/立方毫米 & 500～2000 个/立方毫米 \\
嗜碱性粒细胞绝对数 & 100 个/立方毫米 & $<700$ 个/立方毫米 \\
淋巴细胞绝对数 & 1800 个/立方毫米 & 800～4000 个/立方毫米 \\
单核细胞绝对数 & 400 个/立方毫米 & 100～1000 个/立方毫米 \\
血压（收缩压/舒张压） & 120/75 毫米汞柱 & 90～140/60～90 毫米汞柱 \\
血糖 & 6.5 毫摩尔/升 & 3.9～5.8 毫摩尔/升 \\
\hline
\end{tblr}
\end{center}

人的心脏每跳动一次大约输送 $8\times10^{-5}\Umc$ 血液，根据这位病人体检表的数据，计算他的心脏每收缩一次所做的功大约是 \tkanswer[1.8cm]{1.3}$\UJ$（汞的密度是 $13.6\times10^{3}\Ukgmc$，$g$ 取 $10\Umsq$）

%% answer: 
\jdanswer{
1.3
}

%% memo:
\memoanswer{
本题原卷及材料中未提供详解，待补充。
}

\end{enumerate}

\end{document}
